% ============================================================================
%  Guía: cómo editar (o ampliar) un endpoint en este proyecto To-Do.
%
%  Compilar con:   pdflatex editar-endpoint.tex
%  (No usa babel-spanish para no depender de paquetes extra; los acentos
%   funcionan con inputenc/fontenc + lmodern.)
% ============================================================================
\documentclass[11pt,a4paper]{article}

\usepackage[T1]{fontenc}
\usepackage[utf8]{inputenc}
\usepackage{lmodern}
\usepackage{amssymb}          % \square para la lista de comprobación
\usepackage[margin=2.3cm]{geometry}
\usepackage{xcolor}
\usepackage{listings}
\usepackage{enumitem}
\usepackage[hidelinks]{hyperref}

% ---------- Colores y estilo para los bloques de código ----------
\definecolor{codebg}{RGB}{245,245,247}
\definecolor{codekw}{RGB}{0,90,160}
\definecolor{codecmt}{RGB}{120,120,120}
\definecolor{codestr}{RGB}{160,60,20}
\definecolor{rulecol}{RGB}{210,210,215}

\lstdefinestyle{codigo}{
  backgroundcolor=\color{codebg},
  basicstyle=\ttfamily\small,
  keywordstyle=\color{codekw}\bfseries,
  commentstyle=\color{codecmt}\itshape,
  stringstyle=\color{codestr},
  frame=single,
  rulecolor=\color{rulecol},
  breaklines=true,
  showstringspaces=false,
  columns=fullflexible,
  tabsize=2,
  captionpos=b,
  % 'listings' no entiende UTF-8 por sí solo: mapeamos los caracteres
  % acentuados que aparecen dentro de los bloques de código.
  literate=%
    {á}{{\'a}}1 {é}{{\'e}}1 {í}{{\'i}}1 {ó}{{\'o}}1 {ú}{{\'u}}1
    {Á}{{\'A}}1 {É}{{\'E}}1 {Í}{{\'I}}1 {Ó}{{\'O}}1 {Ú}{{\'U}}1
    {ñ}{{\~n}}1 {Ñ}{{\~N}}1 {ü}{{\"u}}1 {Ü}{{\"U}}1
    {¿}{{?`}}1 {¡}{{!`}}1,
}
\lstset{style=codigo}

\newcommand{\archivo}[1]{\texttt{#1}}

\title{\textbf{Cómo editar un endpoint}\\[2pt]
       \large Proyecto To-Do (Spring Boot + Angular)}
\author{Guía técnica interna}
\date{\today}

\begin{document}
\maketitle

\begin{abstract}
\noindent Esta guía explica, paso a paso y en orden, cómo modificar un endpoint
de la API cuando hay que añadir o cambiar un campo (por ejemplo, añadir el campo
\emph{prioridad} a las tareas). Cubre el recorrido completo: base de datos,
entidad, DTOs, servicio, controlador, pruebas y el frontend de Angular.
Al final hay una \textbf{lista de comprobación} y los \textbf{errores más
comunes} (y cómo evitarlos).
\end{abstract}

\tableofcontents
\vspace{1em}

% ============================================================================
\section{Arquitectura por capas}
% ============================================================================
El backend está organizado en capas. Una petición HTTP atraviesa siempre el
mismo camino, y por eso una modificación suele tocar \emph{todas} las capas de
forma coherente:

\begin{center}
\ttfamily
Cliente (Angular) $\rightarrow$ Controller $\rightarrow$ Service $\rightarrow$
Repository (JPA) $\rightarrow$ Base de datos (PostgreSQL)
\end{center}

\begin{itemize}[leftmargin=1.4em]
  \item \textbf{Entidad} (\archivo{Task.java}): mapea la tabla a un objeto Java (JPA/Hibernate).
  \item \textbf{DTOs} (\archivo{TaskRequest}, \archivo{TaskResponse}): datos que
        \emph{entran} y \emph{salen} de la API. Nunca se expone la entidad directamente.
  \item \textbf{Service} (\archivo{TaskService}): lógica de negocio y transacciones.
  \item \textbf{Controller} (\archivo{TaskController}): expone las rutas HTTP y valida la entrada.
  \item \textbf{Migración Flyway}: crea/versiona el esquema de la base de datos.
\end{itemize}

\noindent\textbf{Regla clave:} la entidad, los DTOs y el esquema de la base de
datos deben describir \emph{los mismos} campos. Si añades un campo a la entidad
pero no a la tabla (o al revés), Hibernate falla al arrancar con \texttt{ddl-auto:
validate}.

% ============================================================================
\section{Paso 0: la regla de oro de las migraciones}
% ============================================================================
\begin{quote}
\textbf{Nunca edites una migración de Flyway que ya se aplicó.} Crea siempre una
migración \emph{nueva} con el siguiente número de versión.
\end{quote}

Flyway guarda un \emph{checksum} (huella) de cada script en la tabla
\texttt{flyway\_schema\_history}. Si modificas un archivo ya aplicado, el
checksum local deja de coincidir con el guardado y la aplicación \textbf{no
arranca}, con un error como:

\begin{lstlisting}[language=bash]
FlywayValidateException: Migration checksum mismatch for migration version 2
-> Applied to database : -509265412
-> Resolved locally    : 1048635791
\end{lstlisting}

Por eso, para añadir la columna \texttt{priority} a una tabla que \emph{ya
existe}, lo correcto es crear \archivo{V3\_\_add\_priority\_to\_tasks.sql}:

\begin{lstlisting}[language=SQL,caption={backend/src/main/resources/db/migration/V3\_\_add\_priority\_to\_tasks.sql}]
-- Añade la columna de prioridad a una tabla ya existente.
-- Se usa un DEFAULT temporal para no romper las filas ya guardadas
-- (la columna es NOT NULL); luego se puede quitar el DEFAULT.
ALTER TABLE tasks
    ADD COLUMN priority VARCHAR(20) NOT NULL DEFAULT 'MEDIA';

ALTER TABLE tasks
    ALTER COLUMN priority DROP DEFAULT;
\end{lstlisting}

\noindent Si \emph{ya} editaste una migración aplicada y el proyecto está roto,
tienes dos opciones para recuperarte en un entorno de desarrollo:
\begin{enumerate}[leftmargin=1.6em]
  \item Ejecutar \texttt{flyway repair} para re-sincronizar el checksum (no
        modifica el esquema, sólo el historial), \emph{o}
  \item Recrear la base de datos desde cero para que Flyway vuelva a aplicar
        todas las migraciones. Sólo válido si puedes perder los datos.
\end{enumerate}

% ============================================================================
\section{Pasos en el backend}
% ============================================================================

\subsection{1. La migración de base de datos}
Crea la migración nueva (ver Paso 0). Comprueba que el tipo y las restricciones
(\texttt{NOT NULL}, longitud) coincidan con lo que declararás en la entidad.

\subsection{2. La entidad JPA}
Añade el atributo, su \texttt{getter/setter} y (si aplica) al constructor de
conveniencia. Un \texttt{enum} se guarda como texto con \texttt{@Enumerated}:

\begin{lstlisting}[language=Java,caption={backend/.../task/Task.java}]
/** Prioridad de la tarea. Se guarda como texto (BAJA / MEDIA / ALTA). */
@Enumerated(EnumType.STRING)
@Column(name = "priority", nullable = false, length = 20)
private TaskPriority priority;

public TaskPriority getPriority()            { return priority; }
public void setPriority(TaskPriority p)      { this.priority = p; }
\end{lstlisting}

\subsection{3. El enum (si el campo es un conjunto cerrado de valores)}
\begin{lstlisting}[language=Java,caption={backend/.../task/TaskPriority.java}]
public enum TaskPriority {
    BAJA,
    MEDIA,   // OJO: el valor real es MEDIA (no "NORMAL")
    ALTA
}
\end{lstlisting}
\textbf{Consistencia:} el enum es la \emph{fuente de verdad} de los valores. El
frontend y los comentarios del SQL deben usar exactamente los mismos.

\subsection{4. Los DTOs (entrada y salida)}
En \archivo{TaskRequest} añade el campo con su validación; en \archivo{TaskResponse}
añádelo también y mapéalo desde la entidad.

\begin{lstlisting}[language=Java,caption={TaskRequest.java (entrada) y TaskResponse.java (salida)}]
// --- TaskRequest: lo que envía el cliente ---
@NotNull(message = "La prioridad es obligatoria")
TaskPriority priority,

// --- TaskResponse: lo que devuelve la API ---
// (añadir el campo al record y al método fromEntity)
task.getPriority(),
\end{lstlisting}

\subsection{5. El servicio}
Propaga el campo tanto al \emph{crear} como al \emph{actualizar}:
\begin{lstlisting}[language=Java,caption={TaskService.java}]
// create(...)
Task task = new Task(request.title(), request.description(),
                     request.status(), request.priority(),
                     request.dueDate(), owner);

// update(...)  -> dirty checking de JPA hace el UPDATE
task.setPriority(request.priority());
\end{lstlisting}

\subsection{6. El controlador}
En este proyecto el \archivo{TaskController} recibe y devuelve DTOs, así que
normalmente \textbf{no hay que cambiar nada}: al ampliar los DTOs, el nuevo campo
viaja solo. Sólo se toca si cambias la \emph{firma} del endpoint (ruta, método
HTTP, parámetros).

\subsection{7. Las pruebas}
Actualiza los tests para que el nuevo campo forme parte de los datos de prueba y
verifica que se refleja en la respuesta:
\begin{lstlisting}[language=Java,caption={TaskServiceTest.java}]
TaskRequest request = new TaskRequest(
    "Comprar pan", "Integral", TaskStatus.PENDIENTE,
    TaskPriority.MEDIA, LocalDate.now());
...
assertThat(result.priority()).isEqualTo(TaskPriority.MEDIA);
\end{lstlisting}

% ============================================================================
\section{Pasos en el frontend (Angular)}
% ============================================================================

\subsection{1. El modelo TypeScript}
Refleja los mismos tipos que el backend en \archivo{core/models/task.model.ts}:
\begin{lstlisting}[language=Java,caption={task.model.ts (el tipo debe coincidir con el enum del backend)}]
export type TaskPriority = 'BAJA' | 'MEDIA' | 'ALTA';

// Añadir 'priority' a las interfaces Task y TaskRequest, y una lista
// reutilizable de opciones con etiqueta legible:
export const TASK_PRIORITY_OPTIONS: { value: TaskPriority; label: string }[] = [
  { value: 'BAJA',  label: 'Baja'  },
  { value: 'MEDIA', label: 'Media' },
  { value: 'ALTA',  label: 'Alta'  },
];
\end{lstlisting}

\subsection{2. El formulario (diálogo de alta/edición)}
Añade el control al \texttt{FormGroup}, la opción al \texttt{<mat-select>} y el
campo tanto al precargar (edición) como al construir el \texttt{TaskRequest}.

\subsection{3. La tabla y los filtros del dashboard}
\textbf{Cuidado con dos trampas frecuentes de Angular Material:}

\begin{enumerate}[leftmargin=1.6em]
  \item Si añades una columna a \texttt{displayedColumns}, \emph{debes} crear su
        \texttt{<ng-container matColumnDef="...">} en la plantilla. Si no, la
        tabla falla en tiempo de ejecución con
        \texttt{"Could not find column with id ..."} (no lo detecta la
        compilación).
  \item Un filtro no funciona sólo con pintar el \texttt{<mat-select>}: hay que
        \emph{suscribirse} a sus \texttt{valueChanges}, usarlo dentro de la
        función de filtrado y resetearlo al limpiar filtros.
\end{enumerate}

% ============================================================================
\section{Verificación}
% ============================================================================
\begin{lstlisting}[language=bash]
# Backend: compilar + pruebas
cd backend && ./mvnw test
# Backend: arrancar (aplica migraciones Flyway al inicio)
./mvnw spring-boot:run

# Frontend: compilar (detecta errores de tipos y de plantilla)
cd frontend && npm run build
# Frontend: servir en desarrollo
npm start
\end{lstlisting}

% ============================================================================
\section{Lista de comprobación}
% ============================================================================
\begin{enumerate}[leftmargin=1.6em]
  \item[$\square$] Migración \textbf{nueva} (no editar una ya aplicada).
  \item[$\square$] Entidad: campo + getter/setter + constructor.
  \item[$\square$] Enum creado (si aplica) y usado como fuente de verdad.
  \item[$\square$] DTOs de entrada y salida actualizados (con validación).
  \item[$\square$] Servicio: crear \emph{y} actualizar.
  \item[$\square$] Pruebas actualizadas y en verde.
  \item[$\square$] Modelo TS con el mismo tipo que el backend.
  \item[$\square$] Formulario: control + opción + carga + envío.
  \item[$\square$] Tabla: \texttt{displayedColumns} \emph{y} \texttt{matColumnDef}.
  \item[$\square$] Filtro suscrito, aplicado y reseteado.
  \item[$\square$] \texttt{./mvnw test} y \texttt{npm run build} sin errores.
\end{enumerate}

% ============================================================================
\section{Errores comunes (y cómo evitarlos)}
% ============================================================================
\begin{description}[leftmargin=0em,style=nextline]
  \item[Checksum de Flyway.] Editar una migración ya aplicada. \emph{Solución:}
        crear una migración nueva; nunca tocar las aplicadas.
  \item[Hibernate \texttt{validate} falla al arrancar.] La entidad y la tabla no
        coinciden. \emph{Solución:} que la columna (tipo, \texttt{NOT NULL},
        longitud) case con la anotación \texttt{@Column}.
  \item[\texttt{Could not find column with id ...}] Añadir la columna a
        \texttt{displayedColumns} sin su \texttt{matColumnDef}. \emph{Solución:}
        crear el \texttt{<ng-container matColumnDef>} correspondiente.
  \item[El filtro/menú "no hace nada".] Un control declarado pero no conectado a
        la lógica (sin \texttt{valueChanges}), o un \texttt{<mat-sidenav>}
        desligado de la señal que lo controla. \emph{Solución:} mantener el
        \texttt{[opened]}/\texttt{valueChanges} enlazados con el estado del
        componente.
  \item[Valores desalineados.] El enum dice \texttt{MEDIA} pero un comentario o
        el frontend dicen \texttt{NORMAL}. \emph{Solución:} el enum manda; todo
        lo demás debe copiar sus valores exactos.
\end{description}

\end{document}
